\begin{proof}
We first prove the necessary inequality for a smooth generator.  Put
\[
 c=g(o),\qquad
 \alpha_i=\partial_rg_{ij}(0,0),\qquad
 \beta_j=\partial_sg_{ij}(0,0).
\]
The common axis restrictions make $\alpha_i$ independent of $j$ and
$\beta_j$ independent of $i$.  Evaluating
\eqref{eq:p071-trivalent-first-balance} and
\eqref{eq:p071-trivalent-second-balance} at the double node gives
\begin{equation}
                         \sum_i\alpha_i=0,
                  \qquad \sum_j\beta_j=0.
 \label{eq:p071-trivalent-linear-balances}
\end{equation}
Joint convexity on the $(i,j)$ branch rectangle supplies the supporting
plane inequality at its corner:
\begin{equation}
 g_{ij}(r,s)\geq c+\alpha_ir+\beta_js.
 \label{eq:p071-trivalent-supporting-plane}
\end{equation}
Multiply \eqref{eq:p071-trivalent-supporting-plane} at $z_a$ by
$\lambda_a$ and sum.  The two moment identities and
\eqref{eq:p071-trivalent-linear-balances} give
\[
 \sum_a\lambda_ag(z_a)
 \geq c
   +M\sum_i\alpha_i+N\sum_j\beta_j
 =c.
\]
This proves \eqref{eq:p071-trivalent-mixed-jensen} at stage zero.

We next justify passage through the full generated hierarchy; a bare appeal
to locality would not suffice because the witnessing neighbourhoods could
shrink.  Branch-rectangle convexity is preserved by pointwise decreasing
limits.  It also survives sheafification: a locally convex extended-real
function on a convex rectangle is convex, as follows by restricting to line
segments and using the one-dimensional local-to-global criterion.

Fix one balanced atomic configuration in a product node chart and scale all
of its radii by $t\in[0,1]$.  For a branch-rectangle convex function $v$, put
\begin{equation}
 F_v(t)=\sum_a\lambda_a
   v(p_{i_a}(tr_a),q_{j_a}(ts_a)).
 \label{eq:p071-trivalent-radial-propagation}
\end{equation}
Every summand is extended-real convex in $t$, so $F_v$ is convex, and
$F_v(0)=v(o)$.  If the mixed inequality holds for the scaled configuration
for all sufficiently small positive $t$, then $0$ is a local minimum of
$F_v$.  Convexity makes it a global minimum on $[0,1]$.  Hence the mixed
inequality at the original configuration follows from its validity on an
arbitrarily small node neighbourhood.

This radial propagation removes the shrinking-domain problem.  At a
decreasing-limit step, the induction hypothesis gives the inequality for
every term on the whole displayed node chart; pass to the limit in the
finite weighted sum when the limiting node value is finite, while the case
$v(o)=-\infty$ is automatic.  At a sheafification step, one member of a
local cover contains $o$ and gives the inequality for all sufficiently
small scaled copies of the fixed configuration; apply
\eqref{eq:p071-trivalent-radial-propagation} to propagate it back to the
original scale.  Transfinite induction through the generated hierarchy now
proves the mixed inequality for every member of $\mathcal P$.

We now audit the explicit function.  Since
\begin{equation}
 r+s-3\min\{r,s\}
       =\max\{-2r+s,r-2s\},
 \label{eq:p071-trivalent-diagonal-convexity}
\end{equation}
the diagonal branch functions are maxima of two affine functions and hence
are jointly convex.  The off-diagonal branch functions are affine.  At
$s=0$ every branch function equals $r$, and at $r=0$ every branch function
equals $s$.  They therefore glue continuously with common axis values.

Fix a second branch $j$ and data as in
\eqref{eq:p071-trivalent-first-separate-data}.  On the node axis one has
\[
                         u(\eta_S,q_j(s_0))=s_0.
\]
Hence
\begin{align*}
 &\sum_a\lambda_a
       u(p_{i_a}(r_a),q_j(s_a))-u(\eta_S,q_j(s_0))\\
 &\qquad=
   \sum_a\lambda_ar_a
      -3\sum_{a:i_a=j}\lambda_a\min\{r_a,s_a\}\\
 &\qquad=3M
      -3\sum_{a:i_a=j}\lambda_a\min\{r_a,s_a\}\\
 &\qquad\geq
   3M-3\sum_{a:i_a=j}\lambda_ar_a=0.
\end{align*}
This proves the complete first-star--line family, including unequal radii
and a varying line coordinate.  Interchanging the two factors proves the
second family.

Finally, for the three diagonal atoms of
\eqref{eq:p071-trivalent-obstruction-violation}, both branchwise radial
moments equal $\rho/3$.  Thus
\eqref{eq:p071-trivalent-first-moment} and
\eqref{eq:p071-trivalent-second-moment} hold, whereas
\[
 u_{ii}(\rho,\rho)=2\rho-3\rho=-\rho.
\]
The mixed inequality fails.  The necessary part already proved excludes
$u$ from the full generated sheaf near $o$ and completes the proof.
\end{proof}
